\chapter{Trực Nhật Và Lao Động}
\chapterintro{Lớp học sạch hay không thường phụ thuộc vào số người biết lẩn đúng lúc}{Đây là một chủ đề học đường quá thật để bỏ qua.}

\begin{lucbatbox}{Lục bát: Đến giờ quét lớp thì bận}
Tan giờ ai cũng còn đây
Nhắc đến trực nhật là bay dần đều
Bạn xin xuống dưới văn phòng
Bạn quên áo bạn kêu đau mất rồi
Chổi lau đứng đó bơ vơ
Một hai chiến sĩ cắn răng làm phần
\end{lucbatbox}

\begin{tutuyetbox}{Thất ngôn tứ tuyệt: Cầm chổi như cầm kiếm}
Bạn em quét lớp rất hào hùng
Tay đưa khí thế tựa ra quân
Chỉ tiếc góc bụi cuối phòng ấy
Vẫn còn trụ lại đến vài lần
\end{tutuyetbox}

\begin{ghichu}
Trực nhật là lúc lớp học bộc lộ rõ nhất ai thật sự sống theo tinh thần tập thể và ai chỉ phát biểu rất hay về tập thể.
\end{ghichu}

\ngancach

\begin{lucbatbox}{Lục bát: Lau bảng tượng trưng}
Bạn kia nhận việc lau bảng
Lau ngang hai vệt tưởng rằng xong xuôi
Phấn còn trắng tận góc trời
Như bao lời hứa đầu môi cuối giờ
Nhìn qua ai cũng ngẩn ngơ
Lao động tối giản đang chờ lên ngôi
\end{lucbatbox}

\begin{tutuyetbox}{Thất ngôn tứ tuyệt: Đến lao động rất đẹp}
Bạn tới lao động rất lung linh
Áo quần gọn ghẽ tóc xinh xinh
Làm thì không thấy bao nhiêu cả
Ảnh lưu niệm lại nổi màn hình
\end{tutuyetbox}